\subsubsection{Gauge couplings in four-dimensional Type I string orbifolds --- arXiv:hep-th/9906039}

\textbf{Authors:} I. Antoniadis, C. Bachas, E. Dudas

\paragraph{主な主張}
4次元$Z_N$ type-I orientifoldの$\mathcal{N}=2$および$\mathcal{N}=1$模型でゲージ結合の1-loop thresholdと幾何学的moduli依存性を計算し、$\mathcal{N}=1$ではそのmoduli依存性が$\mathcal{N}=2$ sectorから生じる一方、twisted NS--NS moduliへのtree-level couplingがゲージ結合統一を左右することを示した\cite{Antoniadis:1999TypeIThresholds}。

\paragraph{新規性}
主にheterotic stringで研究されていたmoduli依存thresholdの解析をtype-I orientifoldへ体系的に拡張し、open-stringの1-loop補正とtwisted closed-string moduliへのdisk couplingを同じ枠組みで求めてeffective supergravityと比較した。

\paragraph{位置付け}
Dixon--Kaplunovsky--Louisによるheterotic側のthreshold計算\cite{Dixon:1991Thresholds}に対応するtype-I側の基礎文献であり、後の$T^2$ KK towerからのfield-theory再構成\cite{Ghilencea:2002StringThresholds}を比較する際のstring-theory benchmarkを与える。
